\section{Background}
Obtaining information about a government service often involves more than finding a website. A citizen must identify the correct procedure, understand the required documents and determine which conditions apply to their situation. In the Bangladeshi materials collected for this project, information is distributed across frequently asked questions, citizen charters, application instructions, fee tables and legal documents. Birth registration, passport applications and driving licences illustrate this difficulty: similar words can refer to different stages of a service.

This creates a gap between how information is published and how citizens ask for it. A document may organize requirements by legal provision, while a user asks a short question in everyday Bangla. Keyword search helps locate exact terms, but wording differences can obscure a relevant passage. Conversely, a result containing the right service name may describe the wrong procedure. Instructions for collecting a passport do not necessarily answer a question about applying for one.

Large language models (LLMs) offer a conversational way to explain information. However, an answer generated from model parameters alone need not reflect the supplied government documents or preserve their conditions. Retrieval-Augmented Generation (RAG) addresses this problem by supplying external passages during generation \cite{lewis2020rag}. Source-grounded retrieval makes the evidence inspectable and permits knowledge updates separately from model training. It improves the basis for answering, without guaranteeing that every retrieved passage or generated statement is correct.

In this context, we developed \textit{Civic.ai -- A Citizen-Centric AI-Powered Government Service Assistant}. Its implemented pipeline, referred to as CivicRAG, connects a predominantly Bangla civic-service corpus to a local question-answering interface. The current knowledge base covers birth and death registration, passport services and Bangladesh Road Transport Authority (BRTA) services. Civic.ai provides document-based guidance; it does not submit applications, access personal government records or replace the responsible authority.

\section{Rational of the Study or Motivation}
Our motivation combines an information-access need with a data-engineering problem. Citizens need understandable explanations, whereas published service information must first be collected and organized before a retrieval system can use it. Preparing this knowledge resource is a central contribution of the project.

The team's needs-analysis questionnaire supports this direction. Among valid responses from consenting adults, 55 of 77 submissions (71.4\%) indicated willingness to use the described AI assistant instead of an intermediary. Required-document guidance was selected by 62 of 77 (80.5\%), and payment or appointment information by 58 of 77 (75.3\%). These findings motivated our emphasis on requirements, procedures and source-traceable guidance. They describe interest in a prospective service, not measured adoption of Civic.ai.

Trust also influenced the design. Even under the questionnaire's hypothetical verified-information condition, 42 of 76 respondents (55.3\%) would recheck an answer. Therefore, the interface exposes evidence and available source links rather than asking users to rely on a response alone. Chapter 4 explains the survey processing and item-specific denominators; the detailed findings belong to the results chapter.

We selected local inference to support experimentation without sending each question and its evidence to a hosted generation service. This creates a trade-off between accessible deployment and response time. We consequently emphasize a manageable three-domain scope, reusable evaluation records and clear evidence boundaries rather than attempting every government service at once.

\section{Problem Statement}
The problem addressed by this thesis is how to turn heterogeneous civic-service documents into a usable, source-traceable Bangla question-answering system. The difficulty begins with the data: headings, conditions, tables and provenance can be lost during conversion or chunking. If these relationships are broken, later retrieval cannot reliably recover the intended meaning.

The next difficulty is applicability. A useful result must concern both the service and the action requested by the user. Finally, the generator must express that evidence correctly: it may otherwise omit a condition, combine unrelated requirements or return an incomplete answer. These stages are connected, but require separate examination.

We therefore investigate three questions:
\begin{enumerate}
\item How can collected materials become reusable domain-specific evidence units while preserving source relationships and unresolved quality issues?
\item How does CivicRAG's hybrid retrieval and evidence-selection configuration compare with a dense-only Simple RAG baseline under matched generation conditions?
\item How can retrieval behaviour and answer quality be evaluated separately to identify the strengths and limitations of the implemented system?
\end{enumerate}

\section{Objective}
Our \textbf{research objectives} are to develop and document the three-domain civic-service dataset, examine the selected RAG configuration against a matched baseline, and analyse retrieval and generation using reproducible evidence. Dataset preparation and representation are described in Chapter 4; comparative observations and evaluation limitations are addressed in Chapter 5.

Our \textbf{system objectives} are to support Bangla questions within a selected service domain, retrieve evidence through semantic and lexical matching, generate a Bangla explanation from the supplied material, and retain the passages and identifiers needed to inspect a response. Additional engineering objectives are local execution, configurable models and reversible changes through version control. These objectives inform the requirements in Chapter 3.

The study adapts established RAG techniques rather than introducing a new foundation model. Indexing documents is knowledge-base construction, not model training. The intended contribution is a developed and evaluated civic-service system supported by a traceable dataset, with claims limited to implemented functions and recorded experiments.

\section{Methodology in Brief }
The methodology separates offline preparation from online answering. Offline, original materials are preserved, prepared text is normalized, and document-aware chunks receive identifiers and metadata. Each chunk retains searchable text and source evidence. BGE-M3 embeds the searchable representation \cite{bgem3}; ChromaDB stores vectors and linked text \cite{chroma}. A separate BM25 index supports lexical retrieval.

Online, the question is normalized and searched within the selected domain. Dense and BM25 rankings are combined using weighted Reciprocal Rank Fusion (RRF), followed by reranking and automatic applicability-based evidence selection. Selected passages are passed to a local generator through Ollama \cite{ollama}. The main paired comparison uses Qwen3:8b \cite{qwen3}, with checks for unsuitable-language output and incomplete generation.

The modular pipeline permits independent inspection of preparation, retrieval and generation. Saved questions, contexts, answers and configuration records support evaluation without repeatedly generating answers. A matched dense-only baseline provides the system comparison, while retrieval-only experiments examine ranked passages separately.

\section{Scopes and Challenges}
\textbf{Current scope.} Civic.ai handles single-turn informational questions using three domain-specific collections. Sources contain Bangla, English and mixed-language content, but the principal interface and comparison target Bangla questions and Bangla responses. The 30-question comparison contains ten passport, ten birth-registration and ten driving-licence questions. Death registration and other BRTA services exist in the corpus but have no dedicated questions in this comparison.

\textbf{Data challenges.} Uneven formatting, scanned text, missing URLs and historical instructions complicate preparation. We retain provenance and review flags rather than treating every converted document as independently verified.

\textbf{Retrieval challenges.} Different procedures share vocabulary, and broad questions can require several sections. Domain selection reduces cross-service interference, but relevance and coverage still require evaluation.

\textbf{Generation challenges.} Local models may omit conditions, combine incompatible evidence or use unnatural Bangla. Output limits can interrupt long checklists. Direct, paraphrased and scenario-based questions use the same pipeline, although the main set does not establish robustness across all these variants.

\textbf{Future expansion.} Additional domains, more varied unseen questions and stronger hardware are possible extensions. Transaction execution, document authentication and live application-status lookup remain outside the implementation. Reused development questions support a bounded comparison, not unrestricted performance claims.

\section{Team Overview}
The project team consists of Md Shihab Sarker, G.Araf Ahmed, Tapu Mohammad Fahad, Fahim Foysal Abir and Mohammad Zunaid. Our work brings together source collection, machine-readable document preparation, pipeline development, experimental review and thesis documentation. These activities are interdependent: preparation changes can affect retrieval, while answer failures can expose missing conditions in evidence.

Source directories, experiment records and Git history keep the outputs traceable. They are used to document the project rather than infer individual responsibilities or hours worked from filenames. Section 3.6 describes the completed development sequence and its controls.

\section{Key Learnings and Insights}
The principal architectural lesson is that finding a related document and answering the question are different tasks. Combining retrieval methods creates opportunities to find evidence, but additional passages can introduce distractions. Selection therefore needs to preserve service conditions rather than simply increase the amount of context.

Dataset development demonstrated the value of original-file preservation, headings and stable identifiers. Bangla processing showed that encoding consistency is necessary but insufficient: normalized text may still contain OCR errors or ambiguous wording. Both structure and language quality affect downstream behaviour.

Finally, evaluation must distinguish retrieval quality, answer support and usefulness. Source IDs make an answer inspectable, but do not prove each statement. A longer response may cover more cases without being more precise. These lessons shape the requirements and methodology in the following chapters.
